% 人工对照底本逐节点转写；独立TikZ重绘。

\begin{tikzpicture}[x=1cm,y=-1cm,every node/.style={anchor=west,align=left,inner sep=2pt,font=\fontsize{9.5}{12.3}\selectfont},edge/.style={draw=black!65,line width=.45pt}]
\node[text width=1.789cm,draw=black!65] (n1) at (0.000,0.666) {中止的成立\\条件};
\node[text width=2.791cm,draw=black!65] (n2) at (2.420,0.275) {未实行终了的中止};
\node[text width=1.456cm,draw=black!65] (n3) at (5.840,0.275) {自动放弃};
\node[text width=2.457cm,draw=black!65] (n4) at (2.420,1.056) {实行终了的中止};
\node[text width=3.125cm,draw=black!65] (n5) at (5.840,1.056) {自动放弃＋有效防止};
\node[text width=2.123cm,draw=black!65] (n6) at (9.760,1.056) {可能的有效性};
\node[text width=2.123cm,draw=black!65] (n7) at (12.780,1.056) {实际的有效性};
\draw[edge] (n2.east)--(5.690,0.275)--(5.690,0.275)--(n3.west);
\draw[edge] (n1.east)--(2.270,0.666)--(2.270,0.275)--(n2.west);
\draw[edge] (n6.east)--(12.630,1.056)--(12.630,1.056)--(n7.west);
\draw[edge] (n5.east)--(9.610,1.056)--(9.610,1.056)--(n6.west);
\draw[edge] (n4.east)--(5.690,1.056)--(5.690,1.056)--(n5.west);
\draw[edge] (n1.east)--(2.270,0.666)--(2.270,1.056)--(n4.west);
\end{tikzpicture}

